\subsection{Morphisms of Barsotti--Tate groups and proper epimorphisms}

\begin{proposition}[C-RD1: proper epimorphisms]\label{prop:C-RD1}
\status{Proved.}
Let $S$ be a noetherian scheme, let $f:S'\to S$ be proper, and suppose
$\mathcal O_S\to f_*\mathcal O_{S'}$ is injective. Put $S''=S'\times_S S'$.
For any two $p$-divisible groups $G,H$ over $S$, the sequence
\[
 \operatorname{Hom}_S(G,H)\longrightarrow
 \operatorname{Hom}_{S'}(G_{S'},H_{S'})
 \rightrightarrows
 \operatorname{Hom}_{S''}(G_{S''},H_{S''})
\]
is an equalizer. Thus the corollary proposed at 100:25 is true. No hypothesis
that $S$ be reduced, normal, or excellent is needed.
\end{proposition}

\begin{proof}
We first prove the precise rigidity statement needed below. If $R$ is
noetherian, $J\subset R$ is nilpotent, and $G,H$ are $p$-divisible groups over
$R$, the reduction map
\[
 \operatorname{Hom}_R(G,H)\longrightarrow
 \operatorname{Hom}_{R/J}(G_{R/J},H_{R/J})
\]
is injective. When $p$ is nilpotent in $R$, this is Katz's
Lemma~1.1.3(2), whose hypotheses are verified for Barsotti--Tate groups in
his proof of Theorem~1.2.1 \cite{Katz1981SerreTate}. Its elementary reason is
that the kernel of reduction in the target formal Lie group is killed by
a fixed power of $p$, whereas the source is $p$-divisible. For completeness,
if $J^2=0$ that kernel is an additive tangent module, hence is killed by
any power of $p$ killing $R$. For general nilpotent $J$, apply this argument
to its successive powers.

For general noetherian $R$, localize at a maximal ideal $\mathfrak m$.
If $p$ is a unit, the finite levels are etale and the assertion is
infinitesimal rigidity for etale schemes. If $p\in\mathfrak m$, reduce
modulo $\mathfrak m^r$ for every $r\geq1$, and use the case just proved.
At each fixed finite level the two coordinate maps differ by entries in
finite locally free modules over $R_{\mathfrak m}$. All these entries lie
in $\bigcap_r\mathfrak m^rR_{\mathfrak m}=0$. Thus every finite-level map
agrees, proving the stated rigidity.

The desired descent assertion is local on $S$, so write $S=\operatorname{Spec}R$.
For each $n$, the functor of group homomorphisms
$G[p^n]\to H[p^n]$ is represented by an affine finitely presented
$S$-scheme: use the affine space of linear maps between their finite
locally free coordinate algebras, and impose the unit, multiplication,
comultiplication, and counit equations. Restriction to the preceding
level gives transition maps. Hence the functor
$X=\underline{\operatorname{Hom}}_S(G,H)$ is represented by their affine
inverse limit. This argument does not assert that $X$ is of finite type.

A homomorphism over $S'$ satisfying the two pullback equalities is a map
$g':S'\to X$. Set
\[
 R_0=\ker\bigl(\Gamma(S',\mathcal O_{S'})
 \rightrightarrows\Gamma(S'',\mathcal O_{S''})\bigr).
\]
The ring $R_0$ is finite as an $R$-module, because it is an $R$-submodule
of the finite module $\Gamma(S',\mathcal O_{S'})$ and $R$ is noetherian.
Also $R\hookrightarrow R_0$. The map $g'$ factors through
$S_0=\operatorname{Spec}R_0$. The map $S'\to S_0$ is proper and
schematically dominant, hence surjective. Two geometric points of $S_0$
over the same geometric point of $S$ lift to points of $S'$; the equalizer
condition identifies them. Consequently $S_0\to S$ is finite, surjective,
and radicial.

Write $X=\operatorname{Spec}C$. Let $B$ be the image of the resulting
$R$-algebra map $C\to R_0$. It is again a finite $R$-module. We have a
finite radicial surjection $Y=\operatorname{Spec}B\to S$, with
$R\hookrightarrow B$, and a closed immersion $Y\hookrightarrow X$.
We claim $\Omega^1_{B/R}=0$. This is a finite $B$-module. The two maps
\[
 B\longrightarrow B\oplus\Omega^1_{B/R},\qquad
 b\longmapsto(b,0),\quad b\longmapsto(b,db)
\]
induce two $p$-divisible-group homomorphisms that agree modulo the
square-zero ideal $\Omega^1_{B/R}$. The test ring is noetherian, so the
rigidity proved above makes them equal. Since $Y\to X$ is a closed
immersion, the two displayed maps agree, and $db=0$ for every $b$.

A finite radicial $R$-algebra with vanishing differentials has, at each
point of $S$, fibre exactly its residue field. The cokernel of $R\to B$
therefore has zero fibre at every point; Nakayama's lemma shows it is zero.
As the map was injective, $B=R$. Thus $g'$ descends to $S$.
Finally, uniqueness follows at every finite level from
$\mathcal O_S\hookrightarrow f_*\mathcal O_{S'}$, since the coordinate
algebras of $G[p^n]$ are finite locally free. This proves the equalizer.
\end{proof}

\begin{proposition}[C-RD2: a limit to infinitesimal rigidity]\label{prop:C-RD2}
\status{Proved; correction to an intermediate claim.}
The functor $\underline{\operatorname{Hom}}(G,H)$ need not be formally
unramified on arbitrary test rings when the base has mixed characteristic.
Consequently ``formally unramified'' in the reading of 100:25 cannot be
used without specifying the allowed thickenings.
\end{proposition}

\begin{proof}
Let $R=\mathbf Z_p\oplus I$ be the square-zero extension with
$I=\mathbf Q_p/\mathbf Z_p$. For $n\geq1$, put
$x_n=p^{-n}+\mathbf Z_p\in I$ and $z_n=1+x_n\in R^\times$.
Then $z_n^{p^n}=1$ and $z_{n+1}^p=z_n$. Sending $1\in\mathbf Z/p^n$
to $z_n$ defines a compatible, nonzero homomorphism
\[
 (\mathbf Q_p/\mathbf Z_p)_R\longrightarrow\mu_{p^\infty,R}
\]
that reduces to zero modulo $I$. The ideal $I$ is square-zero but is not
a finitely generated $\mathbf Z_p$-module. Hence this counterexample is
consistent with Proposition C-RD1: its proof only uses noetherian
thickenings, and passes to a finite schematic image before testing
differentials.
\end{proof}
